% !TeX program = xelatex
% =============================================================
% 模板 07：MCM/ICM 美国大学生数学建模竞赛论文
% 编译： scripts\build.ps1 en-mcm
%
% 格式依据：COMAP 官方《MCM/ICM Contest Instructions》与历年竞赛题面（2024--2026 沿用）：
%          · 全英文，12 pt 字体，单倍行距，四边 1 英寸页边距
%          · 第一页为 Summary Sheet（含 Team Control Number、Problem Chosen、摘要）
%          · 全文（含 Summary Sheet）不超过 25 页，附录计入页数
%          · 不需要封面、目录、承诺书（2016 年后已取消）
% 偏差说明：COMAP 只提供 Word 版 Summary Sheet 模板，没有 .cls；本模板按上述要求
%          手工实现，Summary Sheet 的表格线框与官方 Word 版略有差异。
% =============================================================
\documentclass[12pt,letterpaper]{article}

\usepackage[margin=1in]{geometry}
\usepackage{fontspec}
\setmainfont{Times New Roman}
\usepackage{amsmath}
% 公式字体与正文统一为 Times 风格（默认 Computer Modern 会与 Times 正文明显不搭）
\usepackage{unicode-math}
\setmathfont{TeX Gyre Termes Math}
\usepackage{booktabs}
\usepackage{graphicx}
\usepackage{caption}
\usepackage{fancyhdr}
\usepackage{setspace}
\usepackage[hidelinks]{hyperref}

\captionsetup{font=small,labelsep=colon}
\captionsetup[table]{position=top}
\setlength{\parindent}{0pt}
\setlength{\parskip}{6pt}

% 页脚居中页码
\pagestyle{fancy}
\fancyhf{}
\fancyfoot[C]{\thepage}
\renewcommand{\headrulewidth}{0pt}

\begin{document}

% ======================= 第 1 页：Summary Sheet =======================
\thispagestyle{empty}

\noindent
\begin{tabular}{|p{0.34\textwidth}|p{0.30\textwidth}|p{0.30\textwidth}|}
\hline
\rule{0pt}{2.2em}Team Control Number & \multicolumn{2}{l|}{\rule{0pt}{2.2em}XXXXXXX} \\
\hline
\rule{0pt}{2.2em}Problem Chosen & \multicolumn{2}{l|}{\rule{0pt}{2.2em}A / B / C / D / E / F} \\
\hline
\multicolumn{3}{|l|}{\rule{0pt}{2.2em}2026 MCM/ICM Summary Sheet} \\
\hline
\end{tabular}

\vspace{1.2em}

\begin{center}
  {\large\bfseries Placeholder Title: A Sample MCM/ICM Paper for Template Testing}
\end{center}

\vspace{0.5em}
\begin{center}
  {\bfseries Summary}
\end{center}

This is the placeholder summary. The summary sheet is the most important page of an
MCM/ICM paper: it must state the problem, your modeling approach, the key numerical
results, and the main conclusion in about 400--600 words, without equations or
citations. Replace this paragraph with your own summary.

\vspace{0.6em}
\noindent\textbf{Keywords:} placeholder; mathematical modeling; MCM/ICM

\newpage

% ======================= 正文 =======================
\section{Introduction}

This section restates the problem in your own words, reviews the relevant
background, and states the overall modeling framework. Cite related work with
\verb|\cite{}|, for example~\cite{smith2023mcm}.

\section{Assumptions and Justifications}

\begin{enumerate}
  \item The system is observed at hourly intervals; sub-hour fluctuations are averaged out.
  \item Travel time between adjacent regions is proportional to distance and independent of traffic.
  \item Demand follows a Poisson process with a time-varying rate.
\end{enumerate}

\section{Model Formulation}

Let $x_i(t)$ be the state of region $i$ at time $t$ and $u_i(t)$ the control input.
The dynamics are
\begin{equation}
  \frac{\mathrm{d}x_i}{\mathrm{d}t}=f_i\bigl(x_i(t),u_i(t)\bigr)+\sum_{j\in\mathcal{N}(i)}a_{ij}\bigl(x_j-x_i\bigr),
  \label{eq:dynamics}
\end{equation}
where $\mathcal{N}(i)$ is the neighbor set and $a_{ij}$ is the coupling strength.
Equation~\eqref{eq:dynamics} is the equation placeholder.

Parameter values are summarized in Table~\ref{tab:params}.

\begin{table}[htbp]
  \centering
  \caption{Placeholder Parameter Values}
  \label{tab:params}
  \begin{tabular}{lcc}
    \toprule
    Parameter & Symbol & Value \\
    \midrule
    Coupling strength & $a_{ij}$ & 0.05 \\
    Demand rate       & $\lambda$ & 12 veh/h \\
    Horizon           & $T$      & 24 h \\
    \bottomrule
  \end{tabular}
\end{table}

\section{Results}

With the parameters in Table~\ref{tab:params}, the placeholder model reduces total
waiting time by 18.7\% compared with the baseline policy, consistent with the
findings in~\cite{smith2023mcm}.

\section{Sensitivity Analysis}

Varying $\lambda$ by $\pm20\%$ changes the objective by less than 6\%, so the
conclusion is robust to demand estimation error.

\section{Strengths and Weaknesses}

The model is simple enough to be solved exactly and transparent enough to be
explained in the summary. Its main weakness is the homogeneous-demand assumption.

% ======================= 参考文献 =======================
\bibliographystyle{unsrt}
\bibliography{refs}

\newpage
\appendix
\section{Appendix: Code}

Placeholder appendix. Memo, code, and data sheets go here; they count toward the
25-page limit.

\end{document}
